\scanpage{185}
mais il est naturel de mettre sur $\pi_0(L)$ la topologie image inverse de $K$. [topologie \inserted{des $\pi_0(L)$} des partitions ouvertes de $L$]. $F$ est dit \uline{connexe} si $F\simeq\mathrm{id}_{\mathrm{Ens}}$, i.e. $F$ est T-représentable par l'espace ponctuel [i.e. par un espace connexe].

Considérons deux \deleted{foncteurs} flèches
\diagram{\begin{tikzcd}[ampersand replacement=\&,column sep=1em]
F^*:\quad F^{\deleted{0}0}\arrow[r,shift left=.5ex]\arrow[r,shift right=.5ex]\&F^{\deleted{1}1}\qquad(F^0,F^1\in\underline{\operatorname{Hom}}(\mathrm{Ens},\mathrm{Ens}))
\end{tikzcd}}
et le \uline{foncteur noyau}
\[
\pi^0(F^*)=F=\operatorname{Ker}(F^0\rightrightarrows F^1)
\]

Si $F^0,F^1$ sont \deleted{\unsure{strictement pro|strictement}}\inserted{T-}représentables, alors les deux flèches proviennent de flèches
\diagram{\begin{tikzcd}[ampersand replacement=\&,column sep=1em]
K_0 \& K_1\arrow[l,shift left=.6ex]\arrow[l,shift right=.6ex]
\end{tikzcd}}
et $F$ sera \unsure{égalt|égalemt|également} T-représentable [N.B. Toute $\varprojlim$ de foncteurs T-représentables est T-repr, en prenant une limite inductive sur les espaces dont ils proviennent --, \deleted{\unsure{et}} et le $\pi_0$ associé]. Un \deleted{\unsure{couple}} $F^{\unsure{\alpha|*}}$ comme dessus est dit \uline{connexe} si son noyau est connexe.

Considérons un ``foncteur $\frac12$cosimplicial'', \unsure{ie|de} \inserted{objet $\frac12$\inserted{co}simplicial de} $\underline{\operatorname{Hom}}(\mathrm{Ens},\mathrm{Ens})$
\diagram{\begin{tikzcd}[ampersand replacement=\&,column sep=1em]
F^0\arrow[r,shift left=.6ex]\arrow[r,shift right=.6ex]\&F^1\arrow[r,shift left=1.1ex]\arrow[r]\arrow[r,shift right=1.1ex]\&F^2
\end{tikzcd}}
On lui associe foncteurs $\pi^0(F^*;E)$ et \deleted{\unsure{et|$\xi$}}, $\pi^1(F^*,G)$, où $E$ ens variable, $G$ groupe variable. Si $F^*$ est \uline{coponctué} i.e. muni de $F^*\to\mathrm{id}_{\mathrm{Ens}}^*$, on définit de $\overline{\sim}$ $\pi^0(F^*,\xi;E)$ et $\pi^1(F^*,\xi;G)$, $\xi$

% archive sheet number: 184
